\section{Method}
\label{sec:method}

We present the material-aware learner as a structurally minimal
extension of the geometry-only port-Hamiltonian formulation. The goal is
not simply to add another risk cost. The goal is to let the sensed scene
change the local optimization problem itself: which objective terms are
active, which barriers constrain motion, and which coarse action degrees
of freedom are exposed to the controller. The extension is governed by a
single principle: each new source of environmental complexity enters
through one additive term in the stored energy, and every other change in
the algorithm — new force channel, new parameter subspace, new gradient
pathway, new update timescale — follows from that addition by structural
necessity. We make the principle precise in
Section~\ref{sec:method:enrichment}; the intervening subsections
construct the context-enriched learner under it.

% ----------------------------------------------------------------------
\subsection{Problem setup: local information, latent risk}
\label{sec:method:setup}
% ----------------------------------------------------------------------

At each time step $t$ the agent observes only local information: a goal
displacement $\vec{c}_g - \vec{c}_t$, locally-sensed obstacle or neighbor
features $\{(C_j, R_j, W_j)\}_{j=1}^{N(t)}$ within sensing radius $\hat
d$, and a $P{\times}P$ context patch $P_t$ centred at the current
position. The patch contains a smoothed risk field
$\tilde r(\cdot)\in[0,1]$ and a signed-distance field $\phi(\cdot)$ to
hard hazards or infeasible regions. Depending on the domain,
$\tilde r$ may encode material risk, traffic risk, or a procedurally
constructed local interaction field; the method only requires that the
field be locally queryable and differentiable after smoothing. The patch
is bilinearly resampled at the agent's current position at every
integration step, so forces and costs reflect the actual path traversed.
All comparisons use the same local-information regime; differences come
from the decision mechanism, not from privileged sensing.

The observation therefore defines a different local optimization
instance at every step. The goal term, clearance barriers,
context-risk preference, hazard barrier, and lateral-motion channel form
a product space of potential actions and penalties, but only some
factors are active in a given context. We use ``action-space adaptation''
in this restricted sense: the learner does not enumerate a discrete set
of actions and prune it; it changes the effective continuous degrees of
freedom by changing which force channels have non-negligible magnitude
under the sensed potential field.

% ----------------------------------------------------------------------
\subsection{The shared $(\Pos,\Mom)$-update skeleton}
\label{sec:method:skeleton}
% ----------------------------------------------------------------------

The geometry-only and material-aware learners share an identical
phase-space update structure. Written side by side, with $\tau$ the
integration step, $\Mass$ the (constant) mass matrix, and $R_\theta$
the learned dissipative port:
\begin{equation}
\renewcommand{\arraystretch}{1.6}
\begin{array}{ll}
\textbf{Geom. scaffold (geometry-only)} & \textbf{Ctx-enriched (material-aware)} \\[2pt]
\Mom_{t+1} = \Mom_t
  - \tau\,\nabla_\Pos \Ham_{\mathrm{geom}}
  - \tau R_\theta(\Mom_t)
&
\Mom_{t+1} = \Mom_t
  - \tau\,\nabla_\Pos \Ham_{\mathrm{geom}}
  \;\boldsymbol{-\,\tau\,\nabla_\Pos \Ham_{\mathrm{ctx}}}
  - \tau R_\theta(\Mom_t) \\[6pt]
\Pos_{t+1} = \Pos_t + \tau \Mass^{-1} \Mom_{t+1}
&
\Pos_{t+1} = \Pos_t + \tau \Mass^{-1} \Mom_{t+1}
\end{array}
\label{eq:method:skeleton}
\end{equation}
The $\Pos$-update is identical. The dissipative port is identical. The
geometric force is identical and is carried forward unchanged. The
\emph{only} dynamical difference is the boldfaced term
$-\tau\,\nabla_\Pos \Ham_{\mathrm{ctx}}$ in the context-enriched field $\Mom$-update.
This is the entire structural content of the paper: a single new term,
and its consequences. The remainder of this section traces those
consequences.

A consequence already visible in Eq.~\eqref{eq:method:skeleton}: new
sensory information enters through $\Mom$, not through an external
critic or replanner. The cotangent variable acts as the internal
evaluator that integrates clearance, material risk, and dissipative
correction into a single force balance. The decision field \emph{is}
the policy.

% ----------------------------------------------------------------------
\subsection{Factored stored energy and force channels}
\label{sec:method:energy}
% ----------------------------------------------------------------------

The context-enriched field stored energy decomposes additively as
\begin{equation}
\Ham_\theta(\Pos, \Mom)
\;=\;
\underbrace{\tfrac{1}{2}\Mom^{\!\top}\Mass^{-1}\Mom}_{\Ham_{\mathrm{kin}}}
\;+\;
\underbrace{\beta\, D_{\mathcal{G}}(\Pos, \vec c_g)
            + \textstyle\sum_j \alpha_j\, b_{\mathrm{IPC}}(d_j)}_{\Ham_{\mathrm{geom}}\;\text{(Geom. scaffold)}}
\;+\;
\underbrace{\lambda_s\,\tilde r(\vec c)
            + \lambda_h\, b_{\mathrm{sp}}\!\bigl(\phi(\vec c)\bigr)}_{\Ham_{\mathrm{ctx}}\;\text{(new in Ctx-enriched)}},
\label{eq:method:energy}
\end{equation}
where $b_{\mathrm{IPC}}$ is the incremental potential contact
barrier~\citep{li2020ipc} on per-obstacle clearance $d_j$, and
$b_{\mathrm{sp}}(\phi) = k^{-1}\log(1 + \exp(k(\hat d_\phi - \phi)))$
is a softplus-relaxed inverse SDF activated only when $\phi < \hat
d_\phi$. The dissipative port $R_\theta(\Pos, \Mom) = \gamma \Mass^{-1}
\Mom$ remains structurally separate from the stored
energy~\citep{vanderschaft2014port}, allowing symplectic and dissipative
parts of the field to be analyzed independently. In our experiments
$\Mass = I$, so $\Mass^{-1}\Mom = \Mom$ in code; we keep $\Mass$
explicit in the formalism for general port-Hamiltonian compatibility.
We use the subscript ``ctx'' to emphasize that the same term supports
material contexts in DFC and interaction contexts in highway-env; the
mechanism is the context-conditioned force, not the source of the risk
field.

\begin{figure}[t]
  \centering
  \includegraphics[width=\linewidth]{Master_s_Thesis/figs/ch4_factored_hamiltonian.png}
  \caption{\textbf{Factored stored energy and induced force channels.}
  \emph{(A)} The context-enriched energy $\Ham_\theta$ decomposes additively into
  kinetic, geometric, and context components; the dissipative port
  $R_\theta$ remains structurally separate. \emph{(B)} The geometric
  energy carries goal-attraction and obstacle-barrier potentials; the
  context energy carries soft-risk and hard-hazard terms.
  \emph{(C)} Each potential induces a force channel: $F_{\mathrm{soft}}
  = -\lambda_s \nabla\tilde r$ deflects the trajectory smoothly toward
  lower-risk lanes; $F_{\mathrm{hard}} = -\lambda_h\,
  b_{\mathrm{sp}}'(\phi)\,\nabla\phi$ acts as a soft barrier that
  strengthens sharply near hazard boundaries while remaining
  differentiable. \emph{(D)} The field shifts the agent toward the
  lower-risk feasible maneuver \emph{when one exists}, and remains
  near-static otherwise — a property that follows directly from the
  fact that the force is the gradient of the energy
  (see selectivity paragraph below). Neither channel exists in
  geometry scaffold, and neither can be obtained by reweighting the geometric
  coefficients.}
  \label{fig:method:factored}
\end{figure}

Differentiating $\Ham_{\mathrm{ctx}}$ yields two qualitatively distinct
force channels:
\begin{equation}
F_{\mathrm{ctx}}(\vec c)
\;=\;
\underbrace{-\lambda_s\,\nabla \tilde r(\vec c)}_{F_{\mathrm{soft}}\,:\;\text{lateral deflection}}
\;+\;
\underbrace{-\lambda_h\, b_{\mathrm{sp}}'(\phi(\vec c))\,\nabla\phi(\vec c)}_{F_{\mathrm{hard}}\,:\;\text{hazard repulsion}},
\qquad
b_{\mathrm{sp}}'(\phi) \;=\; -\sigma\!\bigl(k(\hat d_\phi - \phi)\bigr).
\label{eq:method:force}
\end{equation}
The two channels have different regularity. $F_{\mathrm{soft}}$ is a
smooth gradient field acting wherever $\tilde r$ varies; it expresses
preference among feasible maneuvers. $F_{\mathrm{hard}}$ approaches a
step function as $\phi \to 0^+$; it expresses a near-prohibition on
boundary contact, smoothed enough to remain differentiable. Risk does
not enter the algorithm as a cost evaluated after the fact; it enters
as a force that reshapes the dynamics directly.

\paragraph{Contextual objective and action degrees of freedom.}
Eq.~\eqref{eq:method:energy} also clarifies how the objective and
constraint set change with context. The geometric term supplies a
goal-and-clearance scaffold. The soft context term supplies a
reward-like preference over safer feasible directions. The hard hazard
term supplies a constraint-like barrier whose active set changes with
the local SDF. The force decomposition turns these terms into an
effective action product space: longitudinal acceleration and damping,
lateral deflection, obstacle repulsion, and hazard avoidance. The context-enriched field
learns coefficients for this product space, but the sensed gradients
determine which factors can actually act. Thus moving objects and
changing material patches do not require the learner to search all
actions uniformly; they modulate which coarse degrees of freedom are
available and useful in the current scene.

\paragraph{Selectivity by construction.}
The form of $F_{\mathrm{soft}}$ in Eq.~\eqref{eq:method:force} carries
an immediate consequence that we want to make explicit. The lateral
component of $F_{\mathrm{soft}}$ at a configuration $\vec c$ is
$-\lambda_s$ times the lateral component of $\nabla \tilde r(\vec c)$.
When the local risk field has a meaningful lateral gradient i.e.,
when adjacent feasible directions differ in risk $F_{\mathrm{soft}}$
has a nonzero lateral component and the field deflects the agent toward
the safer direction. When the risk field is laterally flat, or when the
hard-hazard/barrier terms close the neighboring directions, the lateral
degree of freedom remains present in the architecture but contributes
negligible useful force. Selectivity is therefore a property of the
context-conditioned gradient field: learning controls the coefficients,
while sensing determines which force directions can act.

\paragraph{Coefficient prediction.}
The new coefficients $(\lambda_s, \lambda_h)$ are predicted by
sigmoid-bounded heads from a CNN encoding of the risk patch $P_t$
fused with the existing goal-context token; the geometric heads
$(\alpha,\beta,\gamma)$ are unchanged. Freezing the geometry scaffold geometry
heads makes the new coefficients the only route by which context can
alter the field, which is what makes the ablations interpretable.

% ----------------------------------------------------------------------
\subsection{Tail-risk objective: well-posed differentiation through the quantile}
\label{sec:method:cvar}
% ----------------------------------------------------------------------

The episodic cost combines a terminal goal error, true rollout arc
length, cumulative soft-risk exposure, and a hard-hazard step count:
\begin{equation}
J(\theta) \;=\;
w_g\,\|\Pos_T{-}\Pos_g\|^2
\;+\; w_\ell\,\sum_t \|\Pos_{t+1}{-}\Pos_t\|
\;+\; w_r\,\sum_t \tilde r(\Pos_t)\,\|\Pos_{t+1}{-}\Pos_t\|
\;+\; w_h\,\sum_t \mathbf{1}\!\bigl[\phi(\Pos_t) < \epsilon\bigr],
\label{eq:method:J}
\end{equation}
each rollout realized by integrating Eq.~\eqref{eq:method:skeleton} for
$H$ steps. We use the true rollout arc length rather than $\tau H$ as
the length proxy, so that detoured paths are correctly penalized.

\paragraph{Why expected cost is insufficient.}
Catastrophic contact, collision, or route failure is a tail event: the
bulk of rollouts may look acceptable while the failures that determine
deployment safety are rare. Expected-cost training is then dominated by
typical-case path quality and allocates too little gradient to the
rollouts where new context force channels must become active.

\paragraph{Tail-risk objective and well-posed gradient.}
We instead optimize $\cvar{\alpha}[J] = \mathbb{E}[J \mid J \geq
\mathrm{VaR}_\alpha(J)]$, the conditional mean of the worst $1{-}\alpha$
fraction of rollouts. Direct differentiation is awkward because
$\mathrm{VaR}_\alpha(J)$ is a non-smooth functional of a
$\theta$-dependent distribution. We use the Rockafellar–Uryasev
variational form~\citep{rockafellar2000optimization}:
\begin{equation}
\cvar{\alpha}\!\bigl[J(\theta)\bigr]
\;=\;
\min_{\eta \in \mathbb{R}}\;
\bigl\{\,\eta + \tfrac{1}{1-\alpha}\,
\mathbb{E}[(J(\theta) - \eta)_+]\bigr\},
\label{eq:method:cvar_var}
\end{equation}
At the minimizing threshold $\eta^*$, the envelope theorem gives the
gradient with respect to $\theta$ while holding $\eta^*$ fixed. This
justifies the empirical estimator
\begin{equation}
\widehat{\cvar{\alpha}}(\theta)
\;=\;
\hat\eta + \frac{1}{(1-\alpha)\,B}\sum_{i=1}^B \bigl(J^{(i)}(\theta) - \hat\eta\bigr)_{+},
\qquad
\hat\eta = \widehat{Q}_\alpha(\{J^{(i)}\}_{i=1}^B),
\label{eq:method:cvar_emp}
\end{equation}
with $\hat\eta$ \emph{detached from the computation graph}.
Backpropagation produces the Rockafellar–Uryasev sample subgradient:
gradient flows only through the $\lceil(1-\alpha)B\rceil$ rollouts that
exceed the empirical quantile, giving a consistent estimator of the
population subgradient under mild regularity~\citep{hong2014monte}. We use
$\alpha = 0.95$, $B = 64$, concentrating each gradient step on the
worst $\sim 3$ rollouts. The role of CVaR here is therefore not just
a different scalarization of the same loss: it routes gradient
through the new force channel exactly on the rollouts where that
channel matters. On those rollouts, $-\tau\nabla\tilde r$ contributes
a non-zero $\partial \Pos_{t+1}/\partial \lambda_s$ at every step,
driving $\lambda_s$ toward values that bias the field away from
high-risk regions or interactions \emph{exactly where the safety signal lives}.

% ----------------------------------------------------------------------
\subsection{Three loops as three computational jobs}
\label{sec:method:loops}
% ----------------------------------------------------------------------

The new gradient pathway operates at a different timescale than the
gradients in the geometry-only learner, which forces a third update
loop. The context-enriched field's adaptation hierarchy comprises three concurrent jobs
(Figure~\ref{fig:method:loops}). The loops are not three speeds of the
same algorithm: they read different inputs, update different objects,
and solve different optimization subproblems. They cannot be collapsed
because they answer different questions.

\begin{figure}[t]
  \centering
  \includegraphics[width=\linewidth]{Master_s_Thesis/figs/ch4_three_loops.png}
  \caption{\textbf{Three loops as three distinct computational jobs.}
  The \emph{segment loop} performs a local rank-1 secant Gauss–Newton
  correction of the active coefficient vector using only segment-level
  observables and residual mismatch. The \emph{episode loop} performs
  global optimization of the meta-navigator parameters via
  Eq.~\ref{eq:method:cvar_emp} and full-rollout statistics. The
  \emph{curriculum loop} performs a statistical advancement decision
  using passivity and coverage records over many episodes. The loops
  operate at different timescales, but more importantly, they consume
  different information and update different objects.}
  \label{fig:method:loops}
\end{figure}

\textbf{Loop 1 -- segment (within episode, deployment-time).}
A rank-1 secant Gauss–Newton corrector adjusts $(\beta, \gamma)$ and
the top-$K$ nearest $\alpha_j$ between consecutive frames using
realized clearance, goal distance, and speed. The Jacobian estimate is
maintained as $\hat J \leftarrow \mu \hat J + (1{-}\mu)\,
\Delta y\,\Delta \zeta^{\!\top}/\|\Delta \zeta\|^2$ and a damped
least-squares step $\Delta \zeta^* = (\hat J^{\!\top}\hat J + \rho I)^{-1}
\hat J^{\!\top}(y_{\mathrm{tgt}} - y)$ is applied with per-coordinate
learning rates and projection to $\mathbb{R}_{\geq 0}$. This loop runs
at every simulator step and uses no extra rollouts. Like the context
force channel itself, the segment loop is selective: it adjusts
coefficients only when realized error suggests a mismatch, and remains
quiet otherwise.

\textbf{Loop 2 -- episode (training-time, between rollouts).}
Backpropagation through the differentiable rollout integrator against
Eq.~\eqref{eq:method:cvar_emp}, plus auxiliary losses:
imitation to a risk-aware target ($L_{\mathrm{traj}}, L_{\mathrm{vel}}$),
multi-start robustness ($L_{\mathrm{multi}}$,
\citep{ellendula2025grlsnam}), a clearance penalty
($L_{\mathrm{clear}}$), and a $\lambda$-entropy regularizer
($L_\lambda$) preventing the risk heads from saturating the sigmoid.
This loop updates the full parameter vector including $\theta_R$ and
the risk encoder $\omega$. Operationally, it learns a contextual policy
over the force-product space: which reward terms, barriers, and
motion degrees of freedom should become active under the sensed scene,
and how strongly.

\textbf{Loop 3 -- curriculum (training-time, between phases).}
The geometry scaffold $\rightarrow$ context-enriched field transition is realized by backward-compatible
initialization: geometry-backbone weights are loaded from a geometry scaffold
checkpoint, and the risk-head outputs are initialized so that
$\lambda_s, \lambda_h \equiv 0$ meaning $F_{\mathrm{ctx}} \equiv 0$
at the start of context-enriched field training and the closed-loop dynamics are
identical to geometry scaffold's at that initialization. Imitation weights are
reduced to $0.3{\times}$ their geometry-scaffold values to allow the CVaR signal
to reshape the field. We monitor a passivity surrogate
$\mathcal{L}_{\mathrm{pass}}$ on a held-out validation
window and advance the curriculum only when its violation rate stays
within tolerance, a checkable, though not formally certified, signal
that the enrichment has not destabilized the closed-loop dynamics.

Algorithm~\ref{alg:ctx_train} gives the consolidated training
procedure for Loops 2 and 3.

\begin{algorithm}[t]
\caption{Context-aware continual Hamiltonian learning (training; Loops 2 \& 3)}
\label{alg:ctx_train}
\begin{algorithmic}[1]
\Require geometry-scaffold checkpoint $\theta^{(1)} = (\alpha,\beta,\gamma,\theta_R)$;
         dataset $\mathcal{D}$ of short rollouts with local context
         patches; CVaR level $\alpha_c{=}0.95$; batch size $B$;
         horizon $H$; step $\tau$.
\Statex \textbf{Loop 3 — backward-compatible geometry-scaffold $\to$ context-enriched initialization.}
\State $\theta^{(2)} \leftarrow \theta^{(1)}$; initialize risk encoder
        $\omega$ and heads so $\lambda_s, \lambda_h \equiv 0$ at output
       \Comment{$F_{\mathrm{ctx}} \equiv 0$ initially}
\State Reduce imitation weights $(w_{\mathrm{traj}}, w_{\mathrm{vel}})
        \leftarrow 0.3\,(w_{\mathrm{traj}}, w_{\mathrm{vel}})$
\For{epoch $e = 1, \dots, E$}
  \Statex \hspace{0.7em}\textbf{Loop 2 — episode loop (one minibatch).}
  \For{minibatch $\{(\Pos_0^{(i)}, \Mom_0^{(i)}, \Pos_g^{(i)},
                     \mathcal{O}^{(i)}, P^{(i)}, H^{(i)})\}_{i=1}^{B}
        \in \mathcal{D}$}
    \State $(\alpha^{(i)},\beta^{(i)},\gamma^{(i)})
            \leftarrow \phi_{\mathrm{geom}}(\mathcal{O}^{(i)},
                       \Pos_g^{(i)}{-}\Pos_0^{(i)})$
            \Comment{geometry transformer, unchanged from Geom. scaffold}
    \State $z_r^{(i)} \leftarrow \psi_\omega(P^{(i)})$;\;
           $\lambda_s^{(i)},\lambda_h^{(i)} \leftarrow$
           $\lambda^{\max}\sigma(\text{MLP}(z_r^{(i)}, z_g^{(i)}))$
    \For{$t = 0, \dots, H^{(i)}{-}1$ (differentiable rollout)}
      \State $(\tilde r_t,\phi_t,\nabla\tilde r_t,\nabla\phi_t)
              \leftarrow \mathrm{BilinearSample}(P^{(i)},\Pos_t^{(i)})$
              \Comment{resample at \emph{current} $\Pos$, not $\Pos_0$}
      \State $F_{\mathrm{soft}} \leftarrow -\lambda_s^{(i)} \nabla\tilde r_t$;\;
             $F_{\mathrm{hard}} \leftarrow -\lambda_h^{(i)}\,
                                          b_{\mathrm{sp}}'(\phi_t) \nabla\phi_t$
              \Comment{Eq.~\eqref{eq:method:force}}
      \State $\Mom_{t+1}^{(i)} \leftarrow \Mom_t^{(i)}
              + \tau(F_{\mathrm{geom}} + F_{\mathrm{soft}}
                     + F_{\mathrm{hard}})
              - \tau\gamma^{(i)} \Mass^{-1}\Mom_t^{(i)}$
              \Comment{$R_\theta = \gamma\Mass^{-1}\Mom$;
                       $\Mass{=}I$ in implementation}
      \State $\Pos_{t+1}^{(i)} \leftarrow \Pos_t^{(i)} + \tau \Mass^{-1}\Mom_{t+1}^{(i)}$
      \State Accumulate $\mathrm{cum\_risk}, \mathrm{hard\_count}, \mathrm{arc}$
    \EndFor
    \State $J^{(i)} \leftarrow$ Eq.~\eqref{eq:method:J}
    \State $\hat\eta \leftarrow \widehat{Q}_{\alpha_c}(\{J^{(i)}\}_{i=1}^B)$
            \Comment{empirical quantile, \textbf{detach from graph}}
    \State $\widehat{\cvar{\alpha_c}} \leftarrow
           \hat\eta + \tfrac{1}{(1-\alpha_c)B}\sum_i (J^{(i)} - \hat\eta)_+$
            \Comment{Eq.~\eqref{eq:method:cvar_emp}; gradient flows
                     only through tail rollouts}
    \State $\mathcal{L} \leftarrow \widehat{\cvar{\alpha_c}}
                + \sum_{k} w_k L_k$
            \Comment{$L_k \in \{L_{\mathrm{traj}}, L_{\mathrm{vel}},
                                L_{\mathrm{fric}}, L_{\mathrm{clear}},
                                L_{\mathrm{multi}}, L_\lambda\}$}
    \State $\theta^{(2)} \leftarrow \theta^{(2)} - \mathrm{lr}\,
            \nabla_{\theta^{(2)}}\mathcal{L}$,
            $\;\|\nabla\|_2$ clipped to $5.0$
  \EndFor
  \Statex \hspace{0.7em}\textbf{Loop 3 — curriculum advancement check.}
  \State Compute passivity surrogate $\mathcal{L}_{\mathrm{pass}}$ on
         validation window
  \If{violation rate $< \eta_{\mathrm{tol}}$ over last $W$ epochs}
    \State Mark phase as advanced; optionally widen risk-class
           distribution in $\mathcal{D}$
  \EndIf
\EndFor
\State \Return $\theta^{(2)}$
\end{algorithmic}
\end{algorithm}

% ----------------------------------------------------------------------
\subsection{The enrichment principle}
\label{sec:method:enrichment}
% ----------------------------------------------------------------------

Sections~\ref{sec:method:skeleton}–\ref{sec:method:loops} construct the
The context-enriched learner as a sequence of additions. They are not independent.
Each is forced by the previous, and we close the section by stating the
underlying principle.

\begin{proposition}[Learnable energy enrichment]
\label{prop:enrichment}
Within the class of port-Hamiltonian learners considered in this paper,
let $\Ham^{(1)}(\Pos, \Mom; \theta^{(1)})$ be a trained reduced
Hamiltonian and $\mathcal{L}^{(1)}$ its training objective, and let
$\Ham^{(2)} = \Ham^{(1)} + \Ham_{\mathrm{new}}(\Pos; \theta_{\mathrm{new}})$
be the enriched energy with $\Ham_{\mathrm{new}} \in C^1$. Then the
context-enriched learner contains, by construction:
{\rm (i)} a new force channel $F_{\mathrm{new}} = -\nabla_\Pos
\Ham_{\mathrm{new}}$ in the $\Mom$-update;
{\rm (ii)} a new parameter subspace $\theta_{\mathrm{new}} \subset
\theta^{(2)}$;
{\rm (iii)} a new gradient pathway
$\partial \Pos_{t+1}/\partial \theta_{\mathrm{new}}$ in the rollout
sensitivity recursion;
{\rm (iv)} a new objective term whose gradient acts non-trivially on
$\theta_{\mathrm{new}}$, generically requiring a new update loop at a
distinct timescale.
Conversely, within the same framework, any of {\rm (i)–(iv)} introduced
\emph{without} the corresponding additive energy term induces an
inconsistency between the dynamics, the objective, and the gradient
flow.
\end{proposition}

The proposition formalizes what reading the method linearly already
shows: the context-enriched field is a structurally consistent enlargement of geometry scaffold
along all four axes within this framework.
Eqs.~\eqref{eq:method:skeleton}–\eqref{eq:method:cvar_emp} realize
(i)–(iv) for the specific case $\Ham_{\mathrm{new}} =
\Ham_{\mathrm{ctx}}$. The proof is direct: differentiating the enriched
Hamiltonian gives the new force term; making that term learnable
introduces the corresponding parameters; unrolling the shared update
creates the sensitivity pathway; and optimizing any loss that depends on
the resulting trajectory sends gradient through that pathway. Introducing
only one of these pieces breaks the equality between the energy, the
dynamics, and the optimized objective.

The practical content of the principle is twofold.
\emph{Positively}, it prescribes how to grow the learner when a new
source of complexity appears (semantic risk, multi-agent coupling,
deformability): introduce one energy term and let the other three
enlargements follow. \emph{Negatively}, it identifies ill-formed
extensions: adding a CVaR objective without a corresponding force
channel, or a force channel without a coefficient subspace, produces an
algorithm whose dynamics and gradient flow disagree on what is being
optimized. Crucially, because the new force channel inherits the
structure of $\nabla \Ham_{\mathrm{new}}$ including where that
gradient is large and where it vanishes its selectivity is
determined by the geometry of the field rather than by training-time
tuning. This is what enables the context-sensitive behavior the paper
characterizes empirically, and what distinguishes this method from
risk-weighted extensions that alter the objective without reshaping
the decision field. Appendix~\ref{app:theory_empirical_checks} lists the
falsifiable rollout metrics used to test these consequences, including
scaffold preservation, no-hallucinated escape, and risk deflection.
